\documentclass[11pt,a4paper]{ltjsarticle}
\usepackage{amsmath,amssymb}
\usepackage{graphicx}
\usepackage{booktabs}
\usepackage{array}
\usepackage{luatexja-fontspec}
\setmainjfont{HaranoAjiMincho-Regular}
\setsansjfont{HaranoAjiGothic-Medium}
\linespread{1.05}\selectfont

% Recommended PDF filename:
% アルゴリズム設計論_repo3_2611193_中井平蔵.pdf
% Example build command:
% latexmk -lualatex -jobname=アルゴリズム設計論_repo3_2611193_中井平蔵 main.tex

\begin{document}

\begin{center}
{\LARGE \textbf{アルゴリズム設計論 第3回レポート}}
\end{center}
\vspace{1em}

\begin{flushright}
提出日: \today \\
学籍番号: 2611193 \\
氏名: 中井 平蔵
\end{flushright}

\section*{1.}
上側の量子ビットを第1量子ビット $q_1$、下側の量子ビットを第2量子ビット $q_2$ とする。
また、2量子ビットの計算基底を

\[
\{\lvert 00 \rangle, \lvert 01 \rangle, \lvert 10 \rangle, \lvert 11 \rangle\}
\]

の順に並べる。

示したいのは、

\[
(H \otimes H)\,\mathrm{CNOT}_{1 \to 2}\,(H \otimes H) = \mathrm{CNOT}_{2 \to 1}
\]

である。
ここで、第1量子ビット $q_1$ を制御量子ビット、
第2量子ビット $q_2$ を標的量子ビットとし、$q_1 = 1$ のときのみ $q_2$ を反転させる
CNOTゲートを$\mathrm{CNOT}_{1 \to 2}$ と表す。
また、第2量子ビット $q_2$ を制御量子ビット、
第1量子ビット $q_1$ を標的量子ビットとし、$q_2 = 1$ のときのみ $q_1$ を反転させる
CNOTゲートを$\mathrm{CNOT}_{2 \to 1}$ と表す。

Hゲートの行列表現は、

\[
H = \frac{1}{\sqrt{2}}
\begin{pmatrix}
1 & 1 \\
1 & -1
\end{pmatrix}
\]

であるため、テンソル積 $H \otimes H$ は以下のようになる。

\begin{align*}
H \otimes H
&= \left( \frac{1}{\sqrt{2}}
\begin{pmatrix}
1 & 1 \\
1 & -1
\end{pmatrix} \right)
\otimes
\left( \frac{1}{\sqrt{2}}
\begin{pmatrix}
1 & 1 \\
1 & -1
\end{pmatrix} \right) \\
&= \frac{1}{2}
\begin{pmatrix}
1 \cdot
\begin{pmatrix}
1 & 1 \\
1 & -1
\end{pmatrix}
&
1 \cdot
\begin{pmatrix}
1 & 1 \\
1 & -1
\end{pmatrix}
\\
1 \cdot
\begin{pmatrix}
1 & 1 \\
1 & -1
\end{pmatrix}
&
-1 \cdot
\begin{pmatrix}
1 & 1 \\
1 & -1
\end{pmatrix}
\end{pmatrix} \\
&= \frac{1}{2}
\begin{pmatrix}
1 & 1 & 1 & 1 \\
1 & -1 & 1 & -1 \\
1 & 1 & -1 & -1 \\
1 & -1 & -1 & 1
\end{pmatrix}
\end{align*}

また、$\mathrm{CNOT}_{1 \to 2}$ の行列表現は、

\[
\mathrm{CNOT}_{1 \to 2} =
\begin{pmatrix}
1 & 0 & 0 & 0 \\
0 & 1 & 0 & 0 \\
0 & 0 & 0 & 1 \\
0 & 0 & 1 & 0
\end{pmatrix}
\]

である。よって、

\begin{align*}
(H \otimes H)\,\mathrm{CNOT}_{1 \to 2}\,(H \otimes H)
&= \frac{1}{4}
\begin{pmatrix}
1 & 1 & 1 & 1 \\
1 & -1 & 1 & -1 \\
1 & 1 & -1 & -1 \\
1 & -1 & -1 & 1
\end{pmatrix}
\begin{pmatrix}
1 & 0 & 0 & 0 \\
0 & 1 & 0 & 0 \\
0 & 0 & 0 & 1 \\
0 & 0 & 1 & 0
\end{pmatrix}
\begin{pmatrix}
1 & 1 & 1 & 1 \\
1 & -1 & 1 & -1 \\
1 & 1 & -1 & -1 \\
1 & -1 & -1 & 1
\end{pmatrix} \\
&=
\begin{pmatrix}
1 & 0 & 0 & 0 \\
0 & 0 & 0 & 1 \\
0 & 0 & 1 & 0 \\
0 & 1 & 0 & 0
\end{pmatrix}
\end{align*}

この行列を各計算基底に作用させると、
第2量子ビット $q_2$ が0のときには第1量子ビット $q_1$ は変化せず、
$q_2$ が1のときに $q_1$ が反転していることが分かる。

以上より、

\[
(H \otimes H)\,\mathrm{CNOT}_{1 \to 2}\,(H \otimes H) = \mathrm{CNOT}_{2 \to 1}
\]

が成立する。したがって、

\begin{center}
\fbox{\begin{minipage}{0.9\linewidth}
CNOTゲートの前後で2つの量子ビットのそれぞれにHゲートを作用させると、
CNOTゲートの制御量子ビットと標的量子ビットが入れ替わり、制御方向が反転することが
行列表現によって示された。
\end{minipage}}
\end{center}

\section*{2.}
\subsection*{(1)}
Aliceは、量子ビット $A1$, $A2$ からなるEPRペア
\[
\lvert \Phi^+ \rangle_{A1A2}
= \frac{1}{\sqrt{2}}
\left(
\lvert 00 \rangle_{A1A2} + \lvert 11 \rangle_{A1A2}
\right)
\]

を持っている。また、Aliceの量子ビット $A3$ とBobの量子ビット $B1$ も、EPRペア
\[
\lvert \Phi^+ \rangle_{A3B1}
= \frac{1}{\sqrt{2}}
\left(
\lvert 00 \rangle_{A3B1} + \lvert 11 \rangle_{A3B1}
\right)
\]

を形成している。

したがって、4量子ビット全体の初期状態は、
\[
\lvert \Psi_0 \rangle
= \lvert \Phi^+ \rangle_{A1A2} \otimes \lvert \Phi^+ \rangle_{A3B1}
= \frac{1}{2}
\left(
\lvert 0000 \rangle + \lvert 0011 \rangle + \lvert 1100 \rangle + \lvert 1111 \rangle
\right)
\]

となる。ここでは、量子ビットの順番を $A1, A2, A3, B1$ とする。

まず、$A2$ を制御量子ビット、$A3$ を標的量子ビットとするCNOTゲートを作用させる。
\[
\mathrm{CNOT}_{A2 \to A3}
\]

により、
\[
\lvert 0000 \rangle \to \lvert 0000 \rangle,\quad
\lvert 0011 \rangle \to \lvert 0011 \rangle,\quad
\lvert 1100 \rangle \to \lvert 1110 \rangle,\quad
\lvert 1111 \rangle \to \lvert 1101 \rangle
\]

となる。したがって、CNOTゲート作用後の状態は、
\[
\lvert \Psi_1 \rangle
= \frac{1}{2}
\left(
\lvert 0000 \rangle + \lvert 0011 \rangle + \lvert 1110 \rangle + \lvert 1101 \rangle
\right)
\]

である。

次に、量子ビット $A2$ にHゲートを作用させると、

\begin{align*}
\lvert \Psi_2 \rangle
&= (I \otimes H \otimes I \otimes I)\lvert \Psi_1 \rangle \\
&= \frac{1}{2\sqrt{2}}
\left[
\lvert 0000 \rangle + \lvert 0100 \rangle
+ \lvert 0011 \rangle + \lvert 0111 \rangle
+ \lvert 1010 \rangle - \lvert 1110 \rangle
+ \lvert 1001 \rangle - \lvert 1101 \rangle
\right]
\end{align*}

となる。

ここで、測定を行い $A2$, $A3$ のビットが $(m,n)$ $(m,n \in \{0,1\})$ となる場合を考える。

\begin{center}
\begingroup
\renewcommand{\arraystretch}{1.4}
\setlength{\tabcolsep}{14pt}
\begin{tabular}{c c}
\hline
\((m,n)\) & 測定後の状態 \\
\hline
\((0,0)\) & $\frac{1}{2\sqrt{2}}\left(\lvert 0000 \rangle + \lvert 1001 \rangle\right)$ \\
\((0,1)\) & $\frac{1}{2\sqrt{2}}\left(\lvert 0011 \rangle + \lvert 1010 \rangle\right)$ \\
\((1,0)\) & $\frac{1}{2\sqrt{2}}\left(\lvert 0100 \rangle - \lvert 1101 \rangle\right)$ \\
\((1,1)\) & $\frac{1}{2\sqrt{2}}\left(\lvert 0111 \rangle - \lvert 1110 \rangle\right)$ \\
\hline
\end{tabular}
\endgroup
\end{center}

次に、$A2$, $A3$ の測定結果ごとに、残った $A1B1$ 成分を抜き出して正規化すると、
\begin{center}
\begingroup
\renewcommand{\arraystretch}{1.4}
\setlength{\tabcolsep}{16pt}
\begin{tabular}{c c}
\hline
\((m,n)\) & $A1B1$ の状態（正規化後） \\
\hline
\((0,0)\) & $\frac{1}{\sqrt{2}}\left(\lvert 00 \rangle + \lvert 11 \rangle\right)$ \\
\((0,1)\) & $\frac{1}{\sqrt{2}}\left(\lvert 01 \rangle + \lvert 10 \rangle\right)$ \\
\((1,0)\) & $\frac{1}{\sqrt{2}}\left(\lvert 00 \rangle - \lvert 11 \rangle\right)$ \\
\((1,1)\) & $\frac{1}{\sqrt{2}}\left(\lvert 01 \rangle - \lvert 10 \rangle\right)$ \\
\hline
\end{tabular}
\endgroup
\end{center}

したがって、$A1$ は変化せず、$B1$ に Pauli 演算子 $\sigma_{m,n}$ を適用すると、
\begin{center}
\begingroup
\renewcommand{\arraystretch}{1.4}
\setlength{\tabcolsep}{14pt}
\begin{tabular}{c c c}
\hline
\((m,n)\) & $\sigma_{m,n}$ & 最終的な出力 \\
\hline
\((0,0)\) & $I$ & $\frac{1}{\sqrt{2}}\left(\lvert 00 \rangle + \lvert 11 \rangle\right)$ \\
\((0,1)\) & $X$ & $\frac{1}{\sqrt{2}}\left(\lvert 00 \rangle + \lvert 11 \rangle\right)$ \\
\((1,0)\) & $Z$ & $\frac{1}{\sqrt{2}}\left(\lvert 00 \rangle + \lvert 11 \rangle\right)$ \\
\((1,1)\) & $XZ$ & $\frac{1}{\sqrt{2}}\left(\lvert 00 \rangle + \lvert 11 \rangle\right)$ \\
\hline
\end{tabular}
\endgroup
\end{center}

よって、$A1B1$ の出力状態は必ず
\begin{center}
\fbox{$\displaystyle \frac{1}{\sqrt{2}}\left(\lvert 00 \rangle + \lvert 11 \rangle\right)$}
\end{center}

\subsection*{(2)}
\begin{center}
\fbox{\strut エンタングルメントスワッピング\strut}
\end{center}

\end{document}
